\subsection{The planar restricted three-body problem} \vskip-2mm\label{sec:3bodyDisc}\text{ }\\
Next we consider the planar restricted three-body problem described in \cite{Sze67}. This problem models the gravitational motion of three bodies in a plane with a negligible mass for one of the bodies, for example the Earth--Moon--Satellite system. The equations of motion can be expressed as a first order system:
\begin{equation}
\boldsymbol F(\boldsymbol x) 
=
\begin{pmatrix}\dot{x_1} - y_1 \\
\dot{x_2} -y_2\\
\dot{y_1} -\left(x_1+2y_2-\dfrac{\alpha(x_1-\beta)}{((x_1-\beta)^2+x_2^2)^{\frac{3}{2}}}-\dfrac{\beta(x_1+\alpha)}{((x_1+\alpha)^2+x_2^2)^{\frac{3}{2}}}\right)\\
\dot{y_2} -\left(x_2-2y_1-\dfrac{\alpha x_2}{((x_1-\beta)^2+x_2^2)^{\frac{3}{2}}}-\dfrac{\beta x_2}{((x_1+\alpha)^2+x_2^2)^{\frac{3}{2}}}\right)\end{pmatrix}=\boldsymbol{0}, \label{3bodySys}
\end{equation}
where $(x_1,x_2)$ is the position of the satellite relative to the center of mass of the Earth and Moon, $\alpha, \beta$ are relative masses of two bodies such that $\alpha+\beta=1$. It is well-known that  \eqref{3bodySys} has a conserved quantity called Jacobi integral $J$ given by
\[
 J(\boldsymbol x)=\dfrac{x_1^2 + x_2^2- y_1^2-y_2^2}{2} + \dfrac{\alpha}{((x_1-\beta)^2+x_2^2)^{\frac{1}{2}}}+\dfrac{\beta}{((x_1+\alpha)^2+x_2^2)^{\frac{1}{2}}}.
\]Moreover, there exists a canonical transformation which turns \eqref{3bodySys} into a Hamiltonian system with $J$ being the effective Hamiltonian in the new coordinates \cite{Sze67}. We shall work directly with \eqref{3bodySys} to illustrate the application of the multiplier method without the need to make or know the existence of such transformation. Using \eqref{multCond1}, the transpose of the associated $1\times 4$ multiplier matrix $\Lambda$ is given by
\[
\Lambda(\bb x)^T = \begin{pmatrix} x_1-\dfrac{\alpha(x_1-\beta)}{((x_1-\beta)^2+x_2^2)^{\frac{3}{2}}}-\dfrac{\beta(x_1+\alpha)}{((x_1+\alpha)^2+x_2^2)^{\frac{3}{2}}} \\
 x_2-\dfrac{\alpha x_2}{((x_1-\beta)^2+x_2^2)^{\frac{3}{2}}}-\dfrac{\beta x_2}{((x_1+\alpha)^2+x_2^2)^{\frac{3}{2}}} \\
 -y_1 \\
 -y_2
\end{pmatrix}.
\]
Note that $J$ is a linear combination of single variable functions of $x_i^2, y_i^2$ and two variable functions of $x_1, x_2$. Specifically, for any permutation $\sigma \in S_4$, $\bb x^{k+v_{\sigma^{-1}(1)}} = (x_1^k, x_2^s, *, *)$ where $s=k,k+1$ and $\bb x^{k+v_{\sigma^{-1}(2)}} = (x_1^r, x_2^k, *, *)$ where $r=k,k+1$. Thus, the discrete multiplier matrix simplifies to
\[
&\Lambda^\tau(\bb x^{k+1}, \bb x^k) =\\
& 
\begin{pmatrix} 
\frac{\Delta}{\Delta x_1} J(\bb x^{k+v_{\sigma^{-1}(1)}}) \\
\frac{\Delta}{\Delta x_2} J(\bb x^{k+v_{\sigma^{-1}(2)}}) \\
\frac{\Delta}{\Delta y_1} J(\bb x^{k+v_{\sigma^{-1}(3)}}) \\
\frac{\Delta}{\Delta y_2} J(\bb x^{k+v_{\sigma^{-1}(4)}})
 \end{pmatrix} =
 \begin{pmatrix} 
\overline{x}_1+\frac{\Delta}{\Delta x_1} \left(\dfrac{\alpha}{((x_1-\beta)^2+(x_2^s)^2)^{\frac{1}{2}}}+\frac{\beta}{((x_1+\alpha)^2+(x_2^s)^2)^{\frac{1}{2}}} \right)\\
\overline{x}_2+\frac{\Delta}{\Delta x_2} \left(\dfrac{\alpha}{((x_1^r-\beta)^2+x_2^2)^{\frac{1}{2}}}+\frac{\beta}{((x_1^r+\alpha)^2+x_2^2)^{\frac{1}{2}}}\right)\\
-\overline{y}_1 \\
-\overline{y}_2
 \end{pmatrix}
\]
By the reciprocal rule of \ref{rule:reciprocal}, chain rule of \ref{rule:chainScalar}, and rational power rule of \ref{rule:rationalPower}, 
\[
\frac{\Delta}{\Delta z}\left(\frac{1}{\sqrt{z}}\right) = -\frac{1}{\sqrt{z^k}\sqrt{z^{k+1}}(\sqrt{z^{k}}+\sqrt{z^{k+1}})}
\]
Combining with the chain rule of \ref{rule:chainScalar} and polynomial rule of \ref{rule:polynomial} gives
\begin{align*}
\frac{\Delta}{\Delta x_1}\left(\frac{1}{((x_1+\alpha)^2+(x_2^s)^2)^{\frac{1}{2}}}\right) &= -\frac{2(\overline{x}_1+\alpha)}{A^{k,s}A^{k+1,s}(A^{k,s}+A^{k+1,s})},\\
\frac{\Delta}{\Delta x_2}\left(\frac{1}{((x_1^r+\alpha)^2+x_2^2)^{\frac{1}{2}}}\right) &= -\frac{2\overline{x}_2}{A^{r,k}A^{r,k+1}(A^{r,k}+A^{r,k+1})}, \\
\frac{\Delta}{\Delta x_1}\left(\frac{1}{((x_1-\beta)^2+(x_2^s)^2)^{\frac{1}{2}}}\right) &= -\frac{2(\overline{x}_1-\beta)}{B^{k,s}B^{k+1,s}(B^{k,s}+B^{k+1,s})}, \\
\frac{\Delta}{\Delta x_2}\left(\frac{1}{((x_1^r-\beta)^2+x_2^2)^{\frac{1}{2}}}\right) &= -\frac{2\overline{x}_2}{B^{r,k}B^{r,k+1}(B^{r,k}+B^{r,k+1})},
\end{align*} where $A^{r,s} := \sqrt{(x_1^{r}+\alpha)^2+(x_2^{s})^2)}, B^{r,s} := \sqrt{(x_1^{r}-\beta)^2+(x_2^{s})^2}$. 
